To enable replication and meta-analysis, each study must publish (or make available under controlled access) the following artifacts.

\subsection*{A1. Study manifest}
A machine-readable manifest (e.g., YAML/JSON) containing: protocol version, candidate invariant statement (I1--I5), sampling description, and the list of systems with identifiers.

\subsection*{A2. System registry}
For each system: source location, acquisition date, commit hashes, build instructions, and any required credentials (or a description of access constraints).

\subsection*{A3. Raw inputs archive}
An immutable snapshot of the raw inputs, or a reproducible procedure to fetch them exactly (e.g., by commit hash). If redistribution is not permitted, provide cryptographic hashes and a retrieval script.

\subsection*{A4. Extractors and normalizers}
Source code (or container images) for extractors and normalizers, including configuration files and documented assumptions.

\subsection*{A5. Architectural representations}
For each system, the normalized representation in a documented schema (e.g., graph format with node/edge types). Include schema version and validation rules.

\subsection*{A6. Measurement outputs}
Per-system measurement vectors, intermediate outputs, and logs sufficient to diagnose measurement failures.

\subsection*{A7. Classification outputs}
Per-system classification labels with justification pointers (which measurements / thresholds caused the label).

\subsection*{A8. Analysis notebooks and figures}
Scripts/notebooks used to compute aggregates, plots, and statistical tests.

\subsection*{A9. Toolchain lockfile}
A pinned, executable environment definition (e.g., container image digest, Nix, or language-specific lockfiles) to rerun extraction and measurement.

\subsection*{A10. Replication instructions}
A minimal ``one-command'' rerun path when feasible, plus step-by-step instructions and expected outputs.